\documentclass[10pt,a4paper]{article}
\usepackage[T1]{fontenc}
\usepackage[utf8]{inputenc}
\usepackage[czech]{babel}
\usepackage{lmodern}
\usepackage{microtype}
\usepackage{geometry}
\geometry{left=14mm,right=14mm,top=18mm,bottom=18mm}
\usepackage[table]{xcolor}
\usepackage{array,longtable,tabularx,booktabs}
\usepackage{fancyhdr,lastpage}
\usepackage{ragged2e}
\usepackage{hyperref}
\hypersetup{hidelinks,pdftitle={Rozpis výuky matematiky – 2. ročník SPŠ, 102 hodin podle tematických plánů a skutečné praxe}}
\setlength{\parindent}{0pt}
\setlength{\parskip}{5pt}
\setlength{\tabcolsep}{4pt}
\renewcommand{\arraystretch}{1.08}
\setlength{\headheight}{18pt}
\pagestyle{fancy}\fancyhf{}
\fancyhead[L]{\footnotesize Matematika -- rozpis výuky}
\fancyhead[R]{\footnotesize 2. ročník SPŠ}
\fancyfoot[L]{\footnotesize Rozpis hodin podle tematického plánu}
\fancyfoot[R]{\footnotesize \thepage/\pageref{LastPage}}
\renewcommand{\headrulewidth}{0.35pt}
\newcommand{\testr}[1]{\textbf{#1}}
\begin{document}
\begin{center}
{\LARGE\bfseries Rozpis výuky matematiky}\\[3mm]
{\large 2. ročník střední průmyslové školy}\\[3mm]
{\small Přehled plánované výuky podle tematických plánů SP\v{S} a reálných zápisů výuky}
\end{center}
\vspace{3mm}
\begin{tabularx}{\textwidth}{@{}Xrrrr@{}}\toprule
\textbf{Ročník} & \textbf{Hodin} & \textbf{Krátké testy} & \textbf{Souhrnné} & \textbf{Testů celkem}\\\midrule
2. ročník & 102 & 8 & 4 & 12\\\midrule
\textbf{Celkem} & \textbf{102} & \textbf{8} & \textbf{4} & \textbf{12}\\\bottomrule
\end{tabularx}
\vspace{4mm}
\section*{Charakter dokumentu}
\begin{itemize}\setlength{\itemsep}{2pt}
\item Dokument obsahuje rozpis výuky matematiky pro 2. ročník SP\v{S} v celkové dotaci \textbf{102 hodin}.
\item Rozpis je členěn po měsících a čtvrtletích. Každý očíslovaný řádek představuje \textbf{jednu plánovanou vyučovací hodinu}.
\item Zkratka \textbf{TP} odkazuje na odpovídající tematický okruh tematického plánu. U testů je označení uvedeno přímo v náplni hodiny.
\item V každém čtvrtletí jsou zařazeny \textbf{tři testy}: dva průběžné testy (zpravidla 15--20 minut) a jedna souhrnná písemná práce.
\end{itemize}

\section*{Z čeho rozpis vychází}
\begin{itemize}\setlength{\itemsep}{2pt}
\item Hodinová dotace vychází z tematického plánu SP\v{S} pro 2. ročník: \textbf{102 hodin}.
\item Obsah a pořadí témat zohledňují reálné zápisy výuky v dostupných třídách, zejména v 1.A (2024/25), 2.A (2025/26), 3.A (2025/26) a 3.D (2025/26).
\item Historické zápisy slouží jako podklad pro reálné pořadí témat a odhad časové náročnosti; nejde o doslovný přepis třídních knih.
\end{itemize}

\section*{Podkladové dokumenty}
\begin{itemize}\setlength{\itemsep}{2pt}
\item tematické plány SP\v{S},
\item zápisy výuky.
\end{itemize}

\section*{Odlišnosti oproti tematickým plánům}
\begin{itemize}\setlength{\itemsep}{2pt}
\item \textbf{2. ročník:} více prostoru dostávají exponenciální rovnice, substituce a procvičování navazujících algebraických úprav.
\end{itemize}
\clearpage

\section*{2. ročník – exponenciální a logaritmické funkce, goniometrie, geometrie}
\textit{Dotace 102 hodin. Čtvrtletní celky a tři testy v každém čtvrtletí jsou součástí hodinové dotace.}
\vspace{2mm}
\footnotesize
\begin{longtable}{@{}>{\RaggedLeft\arraybackslash}p{9mm}>{\RaggedRight\arraybackslash}p{15mm}>{\RaggedRight\arraybackslash}p{150mm}@{}}
\toprule \textbf{Hod.} & \textbf{TP} & \textbf{Náplň vyučovací hodiny} \\\midrule\endfirsthead
\toprule \textbf{Hod.} & \textbf{TP} & \textbf{Náplň vyučovací hodiny} \\\midrule\endhead
\bottomrule\endfoot
\addlinespace[1mm]\multicolumn{3}{@{}l}{\footnotesize\textbf{1. čtvrtletí} -- 2 průběžné testy + 1 souhrnná práce}\\[0.5mm]
\addlinespace[1.1mm]\multicolumn{2}{@{}l}{\textbf{Září}} & \multicolumn{1}{r@{}}{\textbf{9 hodin}}\\
\midrule
1 & 1.1–1.4 & Úvod do druhého ročníku: pravidla, opakování základních funkcí. \\
2 & 1.1–1.4 & Funkce: definiční obor, obor hodnot a graf z prvního ročníku. \\
3 & 1.1–1.4 & Vlastnosti funkcí: monotonie, prostota a intervaly. \\
4 & 1.1–1.4 & Grafické transformace funkcí: posun a změna měřítka. \\
5 & 1.1–1.4 & Krátké opakování funkcí: samostatná práce s grafem. \\
6 & 1.1–1.4 & Prostá funkce a podmínky existence inverzní funkce. \\
7 & 1.1–1.4 & Inverzní funkce: výměna argumentu a hodnoty; symetrie grafu. \\
8 & 1.1–1.4 & Určení předpisu inverzní funkce a kontrola oborů. \\
9 & T1 & \testr{Průběžný test 1 (15–20 min): Funkce a inverzní funkce. Zbytek hodiny: krátký rozbor a navazující procvičování.} \\
\addlinespace[1.1mm]\multicolumn{2}{@{}l}{\textbf{Říjen}} & \multicolumn{1}{r@{}}{\textbf{12 hodin}}\\
\midrule
10 & 1.1–1.4 & Mocninné funkce: kladné a záporné celočíselné exponenty. \\
11 & 1.5, 1.7–1.8 & Exponenciální funkce: definice, základ, graf a vlastnosti. \\
12 & 1.5, 1.7–1.8 & Srovnání růstu a poklesu; základy větší než jedna a mezi nulou a jednou. \\
13 & 1.5, 1.7–1.8 & Exponenciální rovnice: shodné základy a porovnání exponentů. \\
14 & 1.5, 1.7–1.8 & Úpravy exponentů; převod mocnin na společný základ. \\
15 & 1.5, 1.7–1.8 & Exponenciální rovnice pomocí vytýkání mocniny. \\
16 & 1.5, 1.7–1.8 & Exponenciální rovnice se substitucí t=a\textasciicircum{}x>0: první příklady. \\
17 & 1.5, 1.7–1.8 & Kvadratická rovnice po substituci: návrat k x. \\
18 & 1.5, 1.7–1.8 & Substituce ve složitějších exponenciálních rovnicích. \\
19 & 1.5, 1.7–1.8 & Grafy exponenciálních funkcí: konstrukce a čtení. \\
20 & T2 & \testr{Průběžný test 2 (15–20 min): Exponenciální funkce a grafy. Zbytek hodiny: krátký rozbor a navazující procvičování.} \\
21 & 1.5, 1.7–1.8 & Exponenciální rovnice se součty mocnin; účelná úprava. \\
\addlinespace[1.1mm]\multicolumn{2}{@{}l}{\textbf{Listopad}} & \multicolumn{1}{r@{}}{\textbf{9 hodin}}\\
\midrule
22 & 1.5, 1.7–1.8 & Soustavy exponenciálních rovnic: základní postup. \\
23 & 1.5, 1.7–1.8 & Soustavy exponenciálních rovnic se společným základem. \\
24 & 1.5, 1.7–1.8 & Exponenciální rovnice: procvičování vytýkání a substituce. \\
25 & 1.5, 1.7–1.8 & Exponenciální nerovnice: monotonie podle základu. \\
26 & 1.5, 1.7–1.8 & Exponenciální nerovnice s úpravou společného základu. \\
27 & 1.5, 1.7–1.8 & Soustavy exponenciálních rovnic: opakování a rozbor. \\
28 & 1.5, 1.7–1.8 & Technické aplikace exponenciálního růstu a poklesu. \\
29 & P1 & \testr{Čtvrtletní písemná práce: Funkce a exponenciální rovnice.} \\
30 & 1.5, 1.7–1.8 & Exponenciální rovnice: smíšené příklady před souhrnným testem. \\
\addlinespace[1mm]\multicolumn{3}{@{}l}{\footnotesize\textbf{2. čtvrtletí} -- 2 průběžné testy + 1 souhrnná práce}\\[0.5mm]
31 & 1.5, 1.7–1.8 & Soustavy exponenciálních rovnic: rozšířené procvičení. \\
32 & 1.6, 1.9–1.10 & Logaritmus jako inverzní operace k umocňování; podmínky. \\
33 & 1.6, 1.9–1.10 & Dekadický a přirozený logaritmus; základní hodnoty. \\
\addlinespace[1.1mm]\multicolumn{2}{@{}l}{\textbf{Prosinec}} & \multicolumn{1}{r@{}}{\textbf{9 hodin}}\\
\midrule
34 & 1.6, 1.9–1.10 & Věty o logaritmu součinu, podílu a mocniny. \\
35 & 1.6, 1.9–1.10 & Logaritmická funkce: graf, definiční obor a monotonie. \\
36 & 1.6, 1.9–1.10 & Logaritmické rovnice typu log\_a f(x)=b: podmínky. \\
37 & 1.6, 1.9–1.10 & Logaritmické rovnice se stejnými základy: porovnání argumentů. \\
38 & 1.6, 1.9–1.10 & Rovnice s použitím logaritmických vzorců. \\
39 & 1.6, 1.9–1.10 & Rovnice s více logaritmy a podmínkami pro argumenty. \\
40 & 1.6, 1.9–1.10 & Logaritmické rovnice s kvadratickým argumentem. \\
41 & T3 & \testr{Průběžný test 3 (15–20 min): Logaritmické rovnice. Zbytek hodiny: krátký rozbor a navazující procvičování.} \\
42 & 1.6, 1.9–1.10 & Logaritmické nerovnice: základ a směr nerovnosti. \\
\addlinespace[1.1mm]\multicolumn{2}{@{}l}{\textbf{Leden}} & \multicolumn{1}{r@{}}{\textbf{12 hodin}}\\
\midrule
43 & 1.6, 1.9–1.10 & Logaritmické nerovnice s více podmínkami. \\
44 & 1.6, 1.9–1.10 & Smíšené logaritmické rovnice a nerovnice. \\
45 & 1.6, 1.9–1.10 & Opakování exponenciálních a logaritmických rovnic. \\
46 & 1.6, 1.9–1.10 & Příklady s nepřípustnými kořeny a nezbytností zkoušky. \\
47 & 2.1–2.5 & Orientovaný úhel a jednotková kružnice. \\
48 & 2.1–2.5 & Stupňová a oblouková míra; převody a úhly lišící se o otočky. \\
49 & T4 & \testr{Průběžný test 4 (15–20 min): Exponenciální a logaritmické rovnice, základy goniometrie. Zbytek hodiny: krátký rozbor a navazující procvičování.} \\
50 & 2.1–2.5 & Sinus a kosinus obecného úhlu; znaménka v kvadrantech. \\
51 & 2.1–2.5 & Tangens a kotangens; definiční obor, nuly a znaménka. \\
52 & P2 & \testr{Pololetní písemná práce: Exponenciální a logaritmické rovnice, základy goniometrie.} \\
53 & 2.1–2.5 & Hodnoty goniometrických funkcí významných úhlů. \\
54 & 2.1–2.5 & Tabulka základních hodnot a určení úhlu z funkční hodnoty. \\
\addlinespace[1mm]\multicolumn{3}{@{}l}{\footnotesize\textbf{3. čtvrtletí} -- 2 průběžné testy + 1 souhrnná práce}\\[0.5mm]
\addlinespace[1.1mm]\multicolumn{2}{@{}l}{\textbf{Únor}} & \multicolumn{1}{r@{}}{\textbf{9 hodin}}\\
\midrule
55 & 2.1–2.5 & Graf funkce sinus: perioda, amplituda, obor hodnot. \\
56 & 2.1–2.5 & Graf funkce kosinus a transformace grafů. \\
57 & 2.1–2.5 & Grafy tangens a kotangens; asymptoty. \\
58 & 2.1–2.5 & Základní goniometrické vzorce a jejich význam. \\
59 & 2.1–2.5 & Vzorce pro součet a rozdíl úhlů; jednoduché úpravy. \\
60 & 2.1–2.5 & Vzorce pro dvojnásobný úhel a úpravy výrazů. \\
61 & T5 & \testr{Průběžný test 5 (15–20 min): Hodnoty a grafy goniometrických funkcí. Zbytek hodiny: krátký rozbor a navazující procvičování.} \\
62 & 2.1–2.5 & Goniometrické identity: ověření pomocí známých vztahů. \\
63 & 2.1–2.5 & Rovnice sin x=a, cos x=a: všechna řešení. \\
\addlinespace[1.1mm]\multicolumn{2}{@{}l}{\textbf{Březen}} & \multicolumn{1}{r@{}}{\textbf{12 hodin}}\\
\midrule
64 & 2.1–2.5 & Rovnice tg x=a a cotg x=a: periodické řešení. \\
65 & 2.1–2.5 & Goniometrické rovnice s argumentem kx+b. \\
66 & 2.1–2.5 & Goniometrické rovnice se substitucí a úpravou výrazů. \\
67 & 2.6–2.8 & Sinová věta: vztah mezi stranami a protilehlými úhly. \\
68 & 2.6–2.8 & Kosinová věta: výpočet strany obecného trojúhelníku. \\
69 & 2.6–2.8 & Kosinová věta: výpočet úhlu a kontrola výsledku. \\
70 & 2.6–2.8 & Volba mezi sinovou a kosinovou větou podle zadání. \\
71 & 2.6–2.8 & Obecný trojúhelník: dopočítání všech stran a úhlů. \\
72 & 2.6–2.8 & Obsah trojúhelníku a technické úlohy. \\
73 & 2.6–2.8 & Sinová a kosinová věta: smíšené příklady. \\
74 & 2.6–2.8 & Opakování trigonometrie před testem. \\
75 & T6 & \testr{Průběžný test 6 (15–20 min): Sinová a kosinová věta. Zbytek hodiny: krátký rozbor a navazující procvičování.} \\
\addlinespace[1.1mm]\multicolumn{2}{@{}l}{\textbf{Duben}} & \multicolumn{1}{r@{}}{\textbf{3 hodiny}}\\
\midrule
76 & 3.1–3.9 & Základní geometrické pojmy: body, přímky, úhly, rovnoběžnost. \\
77 & P3 & \testr{Čtvrtletní písemná práce: Goniometrie, trigonometrie a technické výpočty.} \\
78 & 3.1–3.9 & Trojúhelníky: úhly, těžnice, výšky, osy. \\
\addlinespace[1mm]\multicolumn{3}{@{}l}{\footnotesize\textbf{4. čtvrtletí} -- 2 průběžné testy + 1 souhrnná práce}\\[0.5mm]
79 & 3.1–3.9 & Kružnice a kruh: tečny, sečny, vzájemná poloha. \\
80 & 3.1–3.9 & Shodnost trojúhelníků a věty o shodnosti. \\
81 & 3.1–3.9 & Podobnost trojúhelníků a dělení úsečky v poměru. \\
82 & 3.1–3.9 & Stejnolehlost: vlastnosti a jednoduché konstrukce. \\
83 & 3.1–3.9 & Pythagorova věta a její obrácení – procvičení. \\
84 & 3.1–3.9 & Euklidovy věty o výšce a odvěsnách. \\
\addlinespace[1.1mm]\multicolumn{2}{@{}l}{\textbf{Květen}} & \multicolumn{1}{r@{}}{\textbf{9 hodin}}\\
\midrule
85 & 3.1–3.9 & Výpočty v pravoúhlém trojúhelníku; vhodná metoda. \\
86 & 3.1–3.9 & Konstrukce trojúhelníku: náčrt a rozbor. \\
87 & 3.1–3.9 & Konstrukce trojúhelníku z kombinace stran a úhlů. \\
88 & T7 & \testr{Průběžný test 7 (15–20 min): Podobnost, Pythagorova a Euklidovy věty. Zbytek hodiny: krátký rozbor a navazující procvičování.} \\
89 & 3.1–3.9 & Konstrukce podle těžnice a výšky; rozbor podmínek. \\
90 & 3.1–3.9 & Množiny bodů dané vlastnosti: osa úsečky a úhlu. \\
91 & 3.1–3.9 & Konstrukční úlohy s kružnicemi a jejich průsečíky. \\
92 & 3.1–3.9 & Středový a obvodový úhel; vztah ke kružnici. \\
93 & 3.1–3.9 & Obvody a obsahy rovinných útvarů; jednotky. \\
\addlinespace[1.1mm]\multicolumn{2}{@{}l}{\textbf{Červen}} & \multicolumn{1}{r@{}}{\textbf{9 hodin}}\\
\midrule
94 & 3.1–3.9 & Složené planimetrické úlohy a opakování. \\
95 & T8 & \testr{Průběžný test 8 (15–20 min): Konstrukce a planimetrie. Zbytek hodiny: krátký rozbor a navazující procvičování.} \\
96 & 4.1–4.4 & Vzájemná poloha bodů, přímek a rovin v prostoru. \\
97 & 4.1–4.4 & Zobrazení těles a prostorová představivost; řezy. \\
98 & 4.1–4.4 & Povrch a objem hranolu a válce. \\
99 & 4.1–4.4 & Povrch a objem jehlanu, kužele a koule. \\
100 & 4.1–4.4 & Komolá tělesa a jednoduché technické úlohy. \\
101 & P4 & \testr{Pololetní písemná práce: Planimetrie, základy stereometrie; souhrn 2. ročníku.} \\
102 & 4.1–4.4 & Závěrečné opakování geometrie a funkcí. \\
\end{longtable}
\normalsize
\end{document}
